\chapter{استثناءات أخطاء الصفحات}
\label{chap:pgfault}

يقوم المعالج بتوليد \indextext{استثناء خطأ الصفحة (Page Fault Exception)} عندما تحاول إحدى التعليمات الوصول إلى عنوان افتراضي لا تقابله مدخلة صالحة في جدول الصفحات،
أو عندما تنتهك العملية بتات الصلاحيات والحماية بالصفحة.

يحفظ المعالج العنوان الافتراضي الذي تسبب في الاستثناء بداخل السجل الخاص \texttt{stval}،
ويحفظ سبب الاستثناء بداخل السجل \texttt{scause}.

تستغل أنظمة التشغيل الحديثة ونواة \texttt{xv6} استثناءات أخطاء الصفحات لتقديم تقنيات متقدمة لإدارة الذاكرة بكفاءة:
\begin{itemize}
\item \textbf{التخصيص الكسول للذاكرة (Lazy Allocation)}: عدم تخصيص صفحات فيزيائية فور استدعاء \texttt{sbrk}، بل الانتظار حتى تحاول العملية الوصول الفعلي للعنوان، فيحدث خطأ صفحة وتخصص النواة الصفحة حينها.
\item \textbf{النسخ عند الكتابة (Copy-on-Write - COW)}: مشاركة الصفحات الفيزيائية بين الأب والابن بقراءة فقط عند \texttt{fork}، وعندما تحاول إحدى العمليات الكتابة، يتولد خطأ صفحة فتنسخ النواة تلك الصفحة المحددة فقط.
\item \textbf{الذاكرة المضمومة بملف (Memory Mapped Files - mmap)}: ربط كتل الملف بالعناوين الافتراضية للعملية وتحميل الصفحات من القرص تدريجياً عند وقوع خطأ الصفحة.
\end{itemize}
